\documentclass[11pt,a4paper]{amsart}

\usepackage[T1]{fontenc}
\usepackage{lmodern}
\usepackage[protrusion=true,expansion=false]{microtype}
\usepackage{mathtools}
\usepackage{amssymb}
\usepackage{xcolor}
\usepackage{hyperref}

\hypersetup{
  colorlinks=true,
  linkcolor=blue!45!black,
  citecolor=blue!45!black,
  urlcolor=blue!45!black,
  pdftitle={An Obstruction to Prime-Local First-Difference Decompositions of Weil's Quadratic Form},
  pdfauthor={Mayk Loide Baccaro}
}

\newcommand{\R}{\mathbb{R}}
\newcommand{\C}{\mathbb{C}}
\newcommand{\PP}{\mathcal{P}^{\!*}}
\newcommand{\ip}[2]{\left\langle #1,#2\right\rangle}
\newcommand{\norm}[1]{\left\lVert #1\right\rVert_{2}}
\setlength{\emergencystretch}{1em}

\theoremstyle{plain}
\newtheorem{theorem}{Theorem}[section]
\newtheorem{proposition}[theorem]{Proposition}
\newtheorem{lemma}[theorem]{Lemma}
\newtheorem{corollary}[theorem]{Corollary}

\theoremstyle{definition}
\newtheorem{definition}[theorem]{Definition}
\newtheorem{remark}[theorem]{Remark}

\title[Prime-local first-difference obstruction]
{An Obstruction to Prime-Local First-Difference\\
Decompositions of Weil's Quadratic Form}
\author[Mayk Loide Baccaro]{Mayk Loide Baccaro\\
Independent researcher\\
\href{mailto:m.baccaro7663@student.nu.edu}{\texttt{m.baccaro7663@student.nu.edu}}}
\date{18 August 2026}
\subjclass[2020]{Primary 11M26; Secondary 42A82}
\keywords{Weil quadratic form, explicit formula, Riemann zeta function,
prime powers, positive quadratic forms, translation operators}

\begin{document}

\begin{abstract}
Weil's criterion makes positivity of an explicit quadratic form equivalent
to the Riemann hypothesis, which motivates attempts to assemble that form
from manifestly positive prime-local pieces.  We isolate the prime-power term
and test a natural such architecture.  For a prime power
$n$, put $w_n=\Lambda(n)/\sqrt n$, $a_n=\log n$, and let
$T_a\varphi(x)=\varphi(x-a)$.  We prove that the scalar-completed local block
\[
 b_n\norm{\varphi}^{2}-2w_n\operatorname{Re}
 \ip{\varphi}{T_{a_n}\varphi}
\]
is nonnegative on $C_c^\infty(\R)$ exactly when $b_n\geq 2w_n$.
The required scalar budget therefore diverges, since
$\sum_{n\geq2}\Lambda(n)/\sqrt n=\infty$.  For every fixed
nonzero compactly supported test function, the ordinary sum of such positive
blocks diverges to $+\infty$.  Consequently, the actual prime-power term
cannot be assembled as an ordinary convergent sum of individually positive
first-difference blocks with a fixed summable scalar completion.  Finite
truncations and summable synthetic weights remain valid.  The obstruction
comes from infinite-weight assembly; each local block still has an exact
algebraic threshold.
\end{abstract}

\maketitle

\section{The obstruction}

Weil's positivity criterion recasts the Riemann Hypothesis as positivity of a
quadratic functional derived from the explicit formula
\cite{Weil1952,Bombieri2000,Conrey2003,Lagarias2007}.  Its prime-power contribution has a
particularly simple translation structure.  That structure invites a local
sum-of-squares construction: assign one positive first-difference block to
each prime power, then assemble the blocks.

This paper determines exactly when that construction works.  Positivity of
the block attached to a prime power of natural weight $w_n$ forces scalar
compensation $b_n\geq2w_n$.  For the actual Weil weights, the resulting total
compensation is infinite.

Let $C_c^\infty(\R)$ carry the inner product
\[
 \ip{\varphi}{\psi}=\int_{\R}\varphi(x)\overline{\psi(x)}\,dx,
 \qquad T_a\varphi(x)=\varphi(x-a).
\]
For $\varphi\in C_c^\infty(\R)$ define
\[
 \widetilde\varphi(x)=\overline{\varphi(-x)},
 \qquad F_\varphi=\varphi*\widetilde\varphi.
\]
Then
\begin{equation}\label{eq:correlation}
 F_\varphi(a)=\ip{\varphi}{T_a\varphi},
 \qquad F_\varphi(-a)=\overline{F_\varphi(a)}.
\end{equation}
Write $\PP=\{p^k:p\text{ prime},\ k\geq1\}$ and, for $n\in\PP$, set
\[
 w_n=\frac{\Lambda(n)}{\sqrt n}>0,
 \qquad a_n=\log n>0.
\]
Given a real scalar $b_n$, consider
\begin{align}
 Q_n^{(b)}(\varphi)
 &=w_n\norm{\varphi-T_{a_n}\varphi}^{2}
   +(b_n-2w_n)\norm{\varphi}^{2} \label{eq:block-difference}\\
 &=b_n\norm{\varphi}^{2}
   -2w_n\operatorname{Re}F_\varphi(a_n). \label{eq:block-correlation}
\end{align}

\begin{theorem}[Prime-local first-difference obstruction]\label{thm:main}
Let $(b_n)_{n\in\PP}$ be real numbers.
\begin{enumerate}
\item The block $Q_n^{(b)}$ is nonnegative on $C_c^\infty(\R)$ if and only if
      $b_n\geq2w_n$.
\item If every block is nonnegative, then
      \[
      \sum_{n\in\PP}b_n=+\infty.
      \]
\item If every block is nonnegative, then for every nonzero
      $\varphi\in C_c^\infty(\R)$,
      \[
      \sum_{n\in\PP}Q_n^{(b)}(\varphi)=+\infty.
      \]
\end{enumerate}
Consequently, no ordinary convergent sum of individually nonnegative blocks
of the form \eqref{eq:block-difference} can represent the prime-power term
with a fixed summable scalar completion.
\end{theorem}

The three assertions separate the mechanism cleanly.  The first is an exact
local positivity threshold.  The second inserts the arithmetic weights.  The
third uses compact support to eliminate the correlation tail.  The divergent
scalar tail remains and prevents positive assembly.

\section{The exact local threshold}

The threshold follows from the elementary approximate translation invariance
of broad normalized plateaux.  This local device is standard; the arithmetic
content begins when its sharp scalar cost is imposed simultaneously at every
actual prime power.

\begin{proposition}\label{prop:local}
Fix $a\neq0$, $w>0$, and $b\in\R$.  The quadratic form
\[
 Q_{a,w,b}(\varphi)
 =w\norm{\varphi-T_a\varphi}^{2}+(b-2w)\norm{\varphi}^{2}
\]
is nonnegative for every $\varphi\in C_c^\infty(\R)$ if and only if
$b\geq2w$.
\end{proposition}

\begin{proof}
If $b\geq2w$, both terms in the displayed expression are nonnegative.

For the converse, choose a nonzero $\eta\in C_c^\infty(\R)$ and set
\[
 \eta_L(x)=L^{-1/2}\eta(x/L),\qquad L>0.
\]
The normalization gives $\norm{\eta_L}=\norm{\eta}$.  The fundamental theorem
of calculus and a change of variables give the quantitative estimate
\begin{equation}\label{eq:plateau-estimate}
 \norm{\eta_L-T_a\eta_L}
 \leq \frac{|a|}{L}\norm{\eta'}.
\end{equation}
Therefore
\[
 Q_{a,w,b}(\eta_L)
 \longrightarrow (b-2w)\norm{\eta}^{2}
 \qquad (L\to\infty).
\]
Nonnegativity for every $L$ forces $b-2w\geq0$.
\end{proof}

For each $n\in\PP$, apply Proposition~\ref{prop:local} with
$a=a_n$, $w=w_n$, and $b=b_n$.  This proves Theorem~\ref{thm:main}(1).
The proof uses the test space itself: every $\eta_L$ is smooth and compactly
supported, although its support expands with $L$.  No completion of the test
space and no spectral approximation are needed.

\section{The arithmetic budget diverges}

The local threshold becomes a global obstruction because the actual
prime-power weights are not summable.

\begin{proposition}\label{prop:weight-divergence}
One has
\[
 \sum_{n\in\PP}\frac{\Lambda(n)}{\sqrt n}=+\infty.
\]
\end{proposition}

\begin{proof}
Let
\[
 W_N=\sum_{1\leq n\leq N}\frac{\Lambda(n)}{\sqrt n},
 \qquad
 \psi(N)=\sum_{1\leq n\leq N}\Lambda(n).
\]
Since $\Lambda(n)\geq0$ and $\sqrt n\leq\sqrt N$ on this range,
\begin{equation}\label{eq:weighted-chebyshev}
 W_N\geq\frac{\psi(N)}{\sqrt N}.
\end{equation}
We use the elementary Chebyshev bound
\begin{equation}\label{eq:chebyshev-lower}
 \psi(N)\geq N\log 2-\log(N+1).
\end{equation}
For completeness, let $L_N=\operatorname{lcm}(1,\ldots,N)$.  Prime-adic
valuation gives $\binom Nk\mid L_N$ for $0\leq k\leq N$; hence
\[
 2^N=\sum_{k=0}^{N}\binom Nk\leq(N+1)L_N.
\]
Taking logarithms and using $\psi(N)=\log L_N$ proves
\eqref{eq:chebyshev-lower}.  Equations \eqref{eq:weighted-chebyshev} and
\eqref{eq:chebyshev-lower} now give
\[
 W_N\geq \sqrt N\log2-\frac{\log(N+1)}{\sqrt N}.
\]
The right-hand side tends to $+\infty$.  Thus $W_N\to+\infty$.  Finally,
$\Lambda(n)$ vanishes away from prime powers, so these are exactly the
increasing-cutoff partial sums in the proposition.
\end{proof}

If every $Q_n^{(b)}$ is nonnegative, Proposition~\ref{prop:local} yields
$b_n\geq2w_n$.  Proposition~\ref{prop:weight-divergence} therefore proves
Theorem~\ref{thm:main}(2).

\section{Failure of ordinary positive assembly}

Compact support makes the prime-power correlation sum finite, but it makes the
positive-block obstruction more transparent.

\begin{proof}[Proof of Theorem~\ref{thm:main}(3)]
Fix nonzero $\varphi\in C_c^\infty(\R)$.  The autocorrelation
$F_\varphi=\varphi*\widetilde\varphi$ is compactly supported.  Hence there is
$A>0$ such that
\[
 F_\varphi(a_n)=0\qquad\text{whenever }a_n>A.
\]
For all sufficiently large $n\in\PP$, equation
\eqref{eq:block-correlation} becomes
\[
 Q_n^{(b)}(\varphi)=b_n\norm{\varphi}^{2}.
\]
Every term is nonnegative, and the tail sum diverges by part (2).  Thus the
ordinary series diverges to $+\infty$.
\end{proof}

In the Fourier convention
\[
 h(r)=\int_{\R}F_\varphi(u)e^{iru}\,du,
 \qquad
 F_\varphi(u)=\frac{1}{2\pi}\int_{\R}h(r)e^{-iru}\,dr,
\]
one has $h(r)\geq0$ on $\R$.  The prime-power part of Weil's explicit formula
in this normalization is
\begin{align}
 P(\varphi)
 &=-\sum_{n\geq2}\frac{\Lambda(n)}{\sqrt n}
     \bigl(F_\varphi(\log n)+F_\varphi(-\log n)\bigr) \label{eq:prime-pair}\\
 &=-2\sum_{n\in\PP}w_n\operatorname{Re}F_\varphi(a_n). \label{eq:prime-real}
\end{align}
This is the sign and factor-of-two convention used in
\cite{Conrey2003}; Appendix~\ref{app:normalization} records the audit.

Suppose an ordinary blockwise-positive construction tried to realize
\begin{equation}\label{eq:target-assembly}
 P(\varphi)+B\norm{\varphi}^{2}
 =\sum_{n\in\PP}Q_n^{(b)}(\varphi),
 \qquad B=\sum_{n\in\PP}b_n<\infty.
\end{equation}
Equations \eqref{eq:block-correlation} and \eqref{eq:prime-real} give the
formal identity, but Theorem~\ref{thm:main} rules out its positive ordinary
assembly: blockwise positivity forces $B=+\infty$, and the right-hand side
diverges on every nonzero test function.

\begin{corollary}\label{cor:weil-scope}
The actual prime-power contribution to Weil's quadratic form admits no
decomposition of the form \eqref{eq:target-assembly} with all of the following
properties:
\begin{enumerate}
\item one uncoupled block is assigned to each prime power;
\item each block is a scalar-completed first translation difference as in
      \eqref{eq:block-difference};
\item each block is nonnegative on all of $C_c^\infty(\R)$; and
\item the scalar completion is an ordinary summable series independent of the
      test function.
\end{enumerate}
\end{corollary}

Corollary~\ref{cor:weil-scope} resolves this architecture.  Cross-prime
couplings, nonlocal quadratic forms, operator-valued completions,
test-function-dependent counterterms, and independently justified
renormalized sums do not satisfy its hypotheses and require separate
analysis.  The theorem neither proves nor disproves the Riemann Hypothesis.

\section{Controls that isolate the mechanism}

Two controls show exactly where the obstruction enters.

\begin{proposition}[Finite-prefix control]\label{prop:finite-control}
For every finite set $E\subset\PP$, the choice $b_n=2w_n$ gives
\[
 \sum_{n\in E}Q_n^{(b)}(\varphi)
 =\sum_{n\in E}w_n\norm{\varphi-T_{a_n}\varphi}^{2}\geq0.
\]
Thus every finite prefix has a valid positive first-difference decomposition.
\end{proposition}

\begin{proposition}[Summable-weight control]\label{prop:summable-control}
Let $(c_j)_{j\geq1}$ be nonnegative with $\sum_jc_j<\infty$, and let
$(\alpha_j)_{j\geq1}\subset\R$.  Then
\[
 \sum_{j\geq1}c_j\norm{\varphi-T_{\alpha_j}\varphi}^{2}
\]
converges for every $\varphi\in L^2(\R)$ and defines a nonnegative quadratic
form.
\end{proposition}

\begin{proof}
Translation is unitary, so
\[
 c_j\norm{\varphi-T_{\alpha_j}\varphi}^{2}
 \leq4c_j\norm{\varphi}^{2}.
\]
The asserted convergence follows by comparison with $4\norm{\varphi}^2
\sum_jc_j$.
\end{proof}

The finite-prefix control verifies the local algebra against the actual
weights.  The summable-weight control verifies the infinite assembly against a
convergent model.  Together they locate the failure in the conjunction of
actual prime-power weights, per-block positivity, and ordinary uncoupled
summation.

\section{Relation to existing formulations}

Weil's original criterion established the global positivity framework
\cite{Weil1952}.  Burnol's place-by-place expositions emphasize how local
terms enter the explicit formula \cite{Burnol1998,Burnol2000}, while Conrey gives a
convenient normalization of the zeta explicit formula and its positivity
criterion \cite{Conrey2003}; Lagarias relates Li coefficients to values of
Weil's quadratic functional on explicit test functions \cite{Lagarias2007}.
Yoshida studied Hermitian forms attached to zeta
functions \cite{Yoshida1992}.  Bombieri analyzed Weil's quadratic functional,
including compact-support restrictions, finite truncations, and associated
spectral questions \cite{Bombieri2000}.  Recent screw-function formulations
organize several positivity approaches through nonlocal operator structures
\cite{Suzuki2026}.  Finite Galerkin dictionaries and numerical realizations
now study truncated or nonlocal Weil operators directly
\cite{Groskin2026,KimEtAl2026}.

The present result addresses a different, sharply specified question: whether
the prime-power term itself can be assembled from uncoupled, individually
positive first-translation-difference blocks using a fixed summable scalar
budget.  Theorem~\ref{thm:main} answers that question exactly.  Its local
threshold is an elementary approximate-invariance calculation; the theorem's
content is the sharp application of that threshold to the actual prime-power
weights and the resulting infinite-assembly exclusion.  The cited
global, finite-dimensional, and nonlocal formulations remain outside the
excluded class.

\appendix

\section{Normalization audit}\label{app:normalization}

This appendix makes the sign, conjugation, and factor of two independently
checkable.  Starting from the definitions,
\begin{align*}
F_\varphi(a)
&=\int_{\R}\varphi(y)\widetilde\varphi(a-y)\,dy\\
&=\int_{\R}\varphi(y)\overline{\varphi(y-a)}\,dy
=\ip{\varphi}{T_a\varphi}.
\end{align*}
Changing variables and conjugating gives
$F_\varphi(-a)=\overline{F_\varphi(a)}$.  Consequently,
\[
 -w\bigl(F_\varphi(a)+F_\varphi(-a)\bigr)
 =-2w\operatorname{Re}\ip{\varphi}{T_a\varphi}.
\]
Unitarity of $T_a$ gives
\begin{align*}
\norm{\varphi-T_a\varphi}^{2}
&=\norm{\varphi}^{2}+\norm{T_a\varphi}^{2}
  -2\operatorname{Re}\ip{\varphi}{T_a\varphi}\\
&=2\norm{\varphi}^{2}
  -2\operatorname{Re}F_\varphi(a).
\end{align*}
Multiplying by $w$ and adding $(b-2w)\norm{\varphi}^{2}$ yields
\eqref{eq:block-correlation}.  These identities fix every sign and numerical
factor used in the main theorem.

\section{Verification remarks}\label{app:reproducibility}

The proof is symbolic and requires no numerical experiment.  Its two key
calculations are the expansion of the translation-difference norm in
\eqref{eq:block-correlation} and the normalized-dilation estimate
\eqref{eq:plateau-estimate}, which together give the sharp threshold
$b\geq2w$.  Proposition~\ref{prop:weight-divergence} then follows from
the lcm form of Chebyshev's lower bound \eqref{eq:chebyshev-lower}, while compact support
of $F_\varphi$ identifies the tail
$Q_n^{(b)}(\varphi)=b_n\norm{\varphi}^{2}$.  The finite-prefix and
summable-weight controls in Propositions~\ref{prop:finite-control} and
\ref{prop:summable-control} record the adjacent regimes in which the same
algebra does assemble correctly.

\paragraph{Data and code availability.}
The theorem and both controls are proved completely in the text and use no
external data or numerical certificate.  An accompanying Lean~4 development
checks the arithmetic divergence with exact prime-power indexing, the abstract
local threshold and scalar obstruction, and both controls.  Its stated scope
does not include the concrete $C_c^\infty$ dilation argument or the
explicit-formula normalization, which remain proved in the text.

\section*{Assistance disclosure}

AI systems were used extensively in mathematical derivation, Lean proof
development, literature discovery, organization, and typesetting.  The author
remains responsible for every mathematical statement, proof, citation, and
submission decision.

\begingroup\footnotesize
\bibliographystyle{amsplain}
\bibliography{references}
\endgroup

\end{document}
